\chapter{1977年江苏连云港市卷}

\begin{problem}
填充

\end{problem}

\begin{problem}
把指数式 \( 3^{2}=9 \) 改写成对数式 . 

\end{problem}

\begin{problem}
写出下列数列的通项公式： （1）2,4,6,8,10,……, \(a_n=\)\fillinblank{}; （2）\(-1\)，1, 2, \(-1\)，1; ……, \(a_n=\)\fillinblank{}.

\end{problem}

\begin{problem}
以点 \((0,0)\) 为圆心，以 \(4\) 为半径的圆的方程是

\end{problem}

\begin{problem}
计算：

\end{problem}

\begin{problem}
（1）\(1^{0 \cdot 123} \times \left(\frac{3}{2}\right)^{-3} \div \left(-\frac{9}{10}\right)^{0}\); （2）\(\log_{2}781 - 1g100\);

\end{problem}

\begin{problem}
\(\frac{\sqrt[3]{x^{\frac{7}{2}}\sqrt{x^{-1}}}}{\sqrt[3]{x^{-7}}\sqrt[3]{x^{13}}}\)，

\[ 4 _{*} \frac {(\sin 2 2 ^{\circ} 3 0 ^{\prime} + \cos 2 2 ^{\circ} 3 0 ^{\prime}) ^{2}}{2 \cos 0 ^{\circ} - 2 \sin 1 5 ^{\circ} \cos 1 5 ^{\circ} + \sin 2 7 0 ^{\circ}}. \]

\end{problem}

\begin{problem}
对于下列函数，写出自变量 \(x\) 的取值范围： （1）\( y = \frac {6 x - 1}{2 x - 1} \)，（2）\( y = \sqrt {x} + \sqrt {- x} \)，（3）\( y = \sqrt {1 9 \frac {1}{2} (x + 1)} \).

\end{problem}

\begin{problem}
求极限 \( \displaystyle\lim_{n\to\infty}\left[\frac{1+3+5+\cdots+(2n-1)}{n+2}-n\right] \)

\end{problem}

\begin{problem}
（1）\(k\) 为何值时，方程 \(x^{2}-2kx+3=0\) 有一个根是 \(-2\)？（2）试判别一元二次方程 \(ax^{2}+(a-b)x-(a+b)=0\) 的根的情况. 

\end{problem}

\begin{problem}
已知椭圆 \(\frac{x^{2}}{5^{2}}+\frac{y^{2}}{3^{2}}=1\)，（1）求此椭圆的长轴的长、离心率、焦点坐标和准线方程；（2）如图，过椭圆的一个焦点 \(F\)，作 \(FA\) 平行于 \(Y\) 轴交椭圆的上半部于 \(A\) 点，求 \(A\) 点的坐标；（3）求 \(\triangle AOF\) 的面积. 

\end{problem}

\begin{problem}
如图：把小矩形 \(ABCD\) 扩充为大矩形 \(AEFG\)，并使

\(BE+EF+FG+GD=22\)，已知小矩形的宽 \(AB=2\)，小矩形的周长为 \(12\)，问 \(BE\) 为何值时扩充部分的面积扩大？最大面积是多少？（扩充部分为图中阴影部分）

\end{problem}

\begin{problem}
已知 \(AD\) 是 \(\odot O\) 的弦，且 \(\widehat{AD}=8\)，半径 \(OH\perp AD\)，交 \(AD\) 于 \(B\) 点，且 \(HB=2\)，又 \(AC\) 是圆的直径，试求 \(BC\) 的长. 附加题：如图，\(A\) 是双曲线 \(x^{2}-y^{2}=1\) 上的一点，\(A\) 点处的法线交坐标轴于 \(B\)、\(C\) 两点，若 \(|BC|=42\)，试求 \(|OA|\). 提示：设 \(A(x_{0},y_{0})\)，先求出法线斜率，次写出方程，找出 \(B\)、\(C\) 坐标，再用距离公式求出 \(A\) 点坐标. 

附加题
\end{problem}

